
\section{Ranking}
\label{sec:6.6}

\begin{definition} 	\label{defn:6.1}
% BEGIN SUBJECT INDEX: defn:6.1
\index{cycle matrix}
\index{rank!matrix}
\index{rank!cycle matrix}
% END SUBJECT INDEX

Given a vector $v$ and matrix $T$, of rank $d$, define the span matrix, or, \textsl{cycle matrix\/}, 
$$
  \Sigma(v,T)=\begin{pmatrix}
  \hbox{---\quad} v T^{\hphantom 2} \hbox{\quad---} \cr 
  \hbox{---\quad}v T^2 \hbox{\quad---} \cr 
  \hbox{---\quad} \ \ \vdots^{\hphantom 2} \ \  \hbox{\quad---} \cr 
  \hbox{---\quad} v T^d \hbox{\quad---} \cr
  \end{pmatrix}
$$
with rows $\{vT,vT^2,\ldots,vT^d\}$.
\end{definition}

\begin{definition} 	\label{defn:6.2}
% BEGIN SUBJECT INDEX: defn:6.2
\index{ranked transformation}
\index{rank!ranked condition}
% END SUBJECT INDEX

Let $w$ be any word in $\cS$. 
Say that $w$ is \textsl{ranked\/} if the set $\{\pi w,\pi w^2,\ldots,\pi w^d\}$ is linearly independent, where $d={\rm rank}(w)$. 
\end{definition}

In other words, $w$ is ranked if $\rk \Sigma(\pi,w)=\rk w$.

\begin{theorem} 	\label{thm:6.1}
If $w$ is ranked, then the kernel, $\cK$, has rank one.
\end{theorem}

\begin{proof} Let $x_i=\pi w^i$ denote the $i^{\rm th}$ row of $\Sigma(\pi,w)$, the corresponding cycle matrix.
Let $e$ be any idempotent in $\cK$. Then by BM-F, we have 
$$
  x_i ee^*=\frac{1}{r}\,u
$$ 
Since every row of $w$ is a linear combination of the $x_i$, every row of $wee^*$ is constant. Since $w$ is a function,
a row of $wee^*$ is some row of $ee^*$. Since $\pi wee^*=(1/r)\,u$, $wee^*$ is not identically zero. So some row of
$ee^*$ is nonzero and constant. Hence that row consists of all 1's, which means that that vertex is in the same block of a partition of $\cK$
as every other vertex. In other words, $\cK$ is rank one.   \qed   
\end{proof}

Let $R$ have a single limiting cycle, $C$, of length $c$. Let $\eta_R$ be the idempotent for the kernel of the semigroup generated by $R$. Define the \textsl{closure\/} of $R$, $\R=\eta_RR$. So $\R$ has rank $c$.

Let $w=\R$ in the above Theorem.
Let $\Sigma'$ be the result of deleting zero columns from $\Sigma(\pi,\R)$.
The condition that $\R$ be ranked becomes that the determinant of $\Sigma'$ be nonzero. 
$\Sigma'$ has the form $CP$, where $C$ is a circulant and $P$ is a permutation. We thus find, after a suitable renumbering of the first row of $M'$, the required condition is 
that its {\it finite Fourier transform never vanishes\/}. 

\subsection{Ranking with $\boldsymbol\Delta$}
\label{sec:6.6.1}

Instead of working only with words from the semigroup, we can consider as well 
$ \Sigma(\pi,\Delta)$. We say $\Delta$ is ranked if $\rk \Sigma(\pi,\Delta)=\rk \Delta$.

\begin{theorem} 	\label{thm:6.2}
If $\Delta$ is ranked, then the kernel, $\cK$, has rank one. 
\end{theorem}

\begin{proof} 
Corollary \ref{cor:pideltan} says $\pi\Delta N=0$. In fact, Corollary \ref{cor:pideltann}
implies further that $\pi \Delta wN=0$ for all $w\in\cS$. Thus,
$\pi\Delta^i N=0$ for all $i\ge1$. If $\Delta$ is ranked, then $\Delta N=0$, and eq. \eqref{eq:N} reduces to
$$ 
  ANA^*=N
$$
Now Proposition \ref{propIRPER} yields $N=0$, i.e., $N_0=J$, which means the kernel is rank one. \qed
\end{proof}

\subsection{Ranking with $\mathbf{A+\vep\boldsymbol\Delta}$}
\label{sec:6.6.2}

Consider mixtures of $R$ and $B$ other than $50$-$50$. Taking $p,q\ge0$, $p+q=1$, we consider
$pR+qB$. Take $|\vep|\le1$ and set $p=(1+\vep)/2,q=(1-\vep)/2$. Then we set
$$
  A_\vep = pR+qB=A+\vep\Delta
$$
recovering $R$ when $\vep=1$ and $B$ when $\vep=-1$.  We say that $A_\vep$ is ranked if
$\rk \Sigma(\pi,A_\vep)=\rk A_\vep$. Note that $\vep=0$ gives back $A$, which cannot be ranked.
We have
 
\begin{theorem} 	\label{thm:6.3}
For $\vep\ne0$, if $A_\vep$ is ranked, then the kernel, $\cK$, has rank one.
\end{theorem}

\begin{proof} 
Calculate 
$$
  \pi A_\vep N=\pi(A+\vep\Delta)N=\pi N=r'u
$$
Similarly, expanding $A_\vep^i$, for $i>1$, apply $\pi$ on the left, using $\pi A=\pi$.
Except for $A^i$, we get terms of the form $\pi\Delta xN$ 
where $x$ is in the semigroup algebra. Hence, all terms vanish except
$\pi A^iN=\pi N=r'u$. If $A_\vep$ is ranked, then every row of $(A+\vep\Delta)N$ is constant.
I.e., $(A+\vep\Delta)N=bJ$, where $b$ is a diagonal matrix. Let $e_i$ be the standard basis row vector.
So the $i^{\rm th}$ row is
$$ 
  e_i(A+\vep\Delta)N=b_iu
$$
Applying this to $\dg \pi.$, using $N\dg\pi.=r'\dg u.$, yields $r'e_i\dg u.=b_i$. I.e., $b=r'I$ or
$$
  (A+\vep\Delta)N=r'J
$$
Writing $A_\vep$ in terms of $R$ and $B$, the $ij$ entry of $A_\vep N$ is
$$
  \underline{A_\vep N}_{\,ij}=pN_{iR,j}+qN_{iB,j}
$$
These are constant as $j$ varies. Taking $j=iR$, $j=iB$ and recalling the symmetry of $N$, we have
$$ 
  pN_{iR,iB}=qN_{iR,iB}
$$
Hence, if $p\ne q$,  $N_{iR,iB}=0$. So in row $iR$, the entries all vanish. By symmetry,
$N_{iB,iR}=0$ and the entries in row $iB$ vanish. But every vertex $i'$ is $iR$ or $iB$ for some $i$,
so $N$ vanishes identically, i.e., the kernel is rank one. \qed
\end{proof}
